\section{退職所得控除}

\subsection{基本問題}

\begin{fillblank}[年金1 2019 問題1(4)\textcircled{4}]
2019 年 3 月に退職所得控除の計算にかかる勤続年数が 30 年となる者が、退職金 2,000 万円
を受け取ったとき、2019 年分の退職所得は（　　A　　）円となる。なお、この他に2019 年
の退職手当等はないものとする。
\end{fillblank}

\begin{answerbox}
  勤続年数20年超なので、
  \begin{flalign*}
  退職所得控除額 &= 800万円 + (勤続年数 - 20年) \times 70万円  \\
     &= 800万円 + (30年 - 20年) \times 70万円 \\
     &= 1,500万円
  \end{flalign*}
  \begin{flalign*}
  退職所得 &= (退職手当等 - 退職所得控除額) \times 1/2 \\
  &= (2,000万円 - 1,500万円) \times 1/2 \\
  &= 250万円
  \end{flalign*}
\end{answerbox}

\begin{truefalse}[年金1 H28 問題1(6)\textcircled{2}]
  \textbf{下線部の内容の正誤を判定し、誤っている場合は正しい内容に改めなさい。}
\\  

退職所得控除額は、勤続年数が 20 年超の場合、\underline{$400 万円＋40 万円 \times（勤続年数－20 年）$}で
計算される。
\end{truefalse}

\begin{answerbox}
  \textbf{誤}

  $800 万円＋70 万円 \times（勤続年数－20 年）$
\end{answerbox}

\newpage

\subsection{特定役員退職手当等に該当する場合}


\begin{truefalse}[年金1 H28 問題1(6)\textcircled{3}]
  \textbf{下線部の内容の正誤を判定し、誤っている場合は正しい内容に改めなさい。}
\\  

役員等勤続年数が 5 年以下である人が支払いを受ける退職手当等のうち、その役員等勤続年
数に対応する退職を起因として支払いを受けるものについては、\underline{退職手当等の額の 3 分の 2 }が
退職所得の金額となる。
\end{truefalse}

\begin{answerbox}
  \textbf{誤}

  退職手当等の額から退職所得控除額を控除した額
\end{answerbox}

\newpage

\subsection{短期退職手当等に該当する場合}

\begin{fillblank}[年金1 2024 問題1(6)ウ]
2020 年4月1日に入社し、2024年4月10日に退職した者に、退職手当等600万円が支払われ
た。この場合の退職所得の金額は（　　A　　）万円となる。なお、役員等の期間はなく、退職一
時金制度からの支払いでありこれ以外の退職手当はないものとする。
\end{fillblank}

\begin{answerbox}
  勤続年数5年以下の場合は、短期退職手当等に該当。

  300万円を超える部分については1/2課税の対象外となる。

   勤続年数20年未満なので、
  \begin{flalign*}
  退職所得控除額 &= 勤続年数 \times 40万円  \\
     &= 5年 \times 40万円 \\
     &= 200万円
  \end{flalign*}
  $600万円 - 200万円 = 400万円 > 300万円$なので$400万円 - 300万円 = 100万円$の部分は1/2課税の対象外。

  \begin{flalign*}
  退職所得 &= 300万円 × 1/2 + 100万円 \\
    &= 250万円
  \end{flalign*}

\end{answerbox}

\begin{explanationbox}
勤続年数の端数は切り上げ。

（例：勤続期間11年2ヶ月の人は、勤続年数12年として退職所得控除額を計算する）
\end{explanationbox}

\begin{shortanswer}[年金1 2021 問題2(5)（改題）]
\begin{enumerate}[(1)]
  \item 退職手当等が「短期退職手当等」に該当する場合の退職所得の計算方法について簡記しなさい。
  \item 以下のＣさんの退職所得の金額を計算しなさい。
\end{enumerate}

＜Ｃさんの情報＞
\begin{itemize}
  \item 勤続年数5 年0 か月で退職
  \item 1,000 万円の退職金を受け取った（法人役員等ではないものとする）
\end{itemize}

\end{shortanswer}

\begin{answerbox}
  \begin{enumerate}[(1)]
    \item \begin{itemize}
      \item 法人役員等以外についても勤続年数5年以下の短期の退職金については、2分の1課税の平準化措置の適用から除外する。
      \item ただし、雇用の流動化等に配慮し、退職所得控除を除いた支払額300万円までは引き続き2分の1課税の平準化措置を適用する。
    \end{itemize}
    \item 短期退職手当等に該当するので、
  300万円を超える部分については1/2課税の対象外となる。

   勤続年数20年未満なので、
  \begin{flalign*}
  退職所得控除額 &= 勤続年数 \times 40万円  \\
     &= 5年 \times 40万円 \\
     &= 200万円
  \end{flalign*}
  $1,000万円 - 200万円 = 800万円 > 300万円$なので$800万円 - 300万円 = 500万円$の部分は1/2課税の対象外。

  \begin{flalign*}
  退職所得 &= 300万円 × 1/2 + 500万円 \\
    &= 650万円
  \end{flalign*}
  \end{enumerate}
\end{answerbox}

\newpage

\subsection{退職所得時控除額の算定における例外}

\begin{truefalse}[年金1 H28 問題1(6)\textcircled{4}]
    \textbf{下線部の内容の正誤を判定し、誤っている場合は正しい内容に改めなさい。}
\\  

  障害者になったことが退職の直接の原因である場合の退職所得控除額は\underline{ 80 万円}加算される。
\end{truefalse}

\begin{answerbox}
  \textbf{誤}

  100万円
\end{answerbox}

\begin{explanationbox}
  類題：年金1 2020 問題1(6)
\end{explanationbox}

\begin{truefalse}[年金1 H29 問題2(3)\textcircled{1}]
      \textbf{以下の文章の正誤を判定し、誤っている場合は正しい内容に改めなさい。}
\\  

   勤続年数が 1 年のとき退職所得控除額は 40 万円、勤続年数が 2 年 6 か月のとき退職所得控除
額は 120 万円となる。なお、前年以前に支払われた他の退職手当等の勤続期間との重複期間は
ないものとする。
\end{truefalse}

\begin{answerbox}
  \textbf{誤}

  勤続年数2年以下の場合、退職所得控除額は80万円。（×40万円）
\end{answerbox}

\begin{fillblank}[年金1 2020 問題1(6)\textcircled{1}]
  Ａさんは退職手当等の額が 2,000 万円、勤続年数が 30 年0 か月である。このときＡさんの
退職所得控除額は（　　A　　）万円であり、退職所得の額は（　　B　　）万円である。なお、障害
者になったことが退職の直接の原因であるものとし、この他に退職手当等はないものとする。
\end{fillblank}

\begin{answerbox}
  勤続20年超であり、障害者になったことが退職の直接の原因であるので、
  \begin{flalign*}
  退職所得控除額 &= 800万円 + (勤続年数 - 20年) \times 70万円  + 100万円 \\
     &= 800万円 + (30年 - 20年) \times 70万円 + 100万円\\
     &= 1,600万円 \qquad \dots (A)
  \end{flalign*}
  \begin{flalign*}
  退職所得 &= (退職手当等 - 退職所得控除額) \times 1/2 \\
  &= (2,000万円 - 1,600万円) \times 1/2 \\
  &= 200万円 \qquad \dots (B)
  \end{flalign*}
\end{answerbox}

\newpage

\subsection{同じ年に２箇所以上から退職所得等が支払われる場合}

\begin{fillblank}[年金1 2020 問題1(6)\textcircled{2}]
  ＢさんはＸ社とＹ社に重複した就労期間を有しており、就職日（入社日）、退職日その他情報
が下表のとおりであるとき、Ｘ社における退職所得控除額は（　　A　　）万円、Ｙ社における退職
所得控除額は（　　B　　）万円である。なお、下表の他に退職手当等はないものとする。

\begin{table}[htb]
  \centering
  \begin{tabular}{|c|c|c|c|c|}
    \hline
     & 就職日（入社日） & 退職日 & 退職手当等の額 &  支払月 \\ \hline
    X社 & 2011 年 4 月 1 日 & 2020 年3 月 31 日 & 500 万円 & 2020 年4 月 \\ \hline
    Y社 & 2013 年4 月1 日& 2020 年8 月 31 日& 200 万円 & 2020 年 10 月 \\ \hline
  \end{tabular}
\end{table}
\end{fillblank}

\begin{answerbox}
  \begin{enumerate}[A:]
    \item X社の勤続期間は9年なので、$40万円 \times 9年 = 360万円$。
    \item Y社の退職金受取時には、2つの勤続期間のうち長い方の勤続期間に重複していない期間を足した期間を勤続期間とする。計算すると9年5ヶ月となるので、端数は切り上げて、勤続期間は10年。\\
よって退職所得控除額は、$40万円 \times 10年 = 400万円$。
  \end{enumerate}
  
\end{answerbox}

\begin{fillblank}[自作 2026/10/20]
  前問のX社、Y社の退職所得をそれぞれ求めよ。
\end{fillblank}

\begin{answerbox}
  \begin{enumerate}[A:]
    \item X社の退職所得は、$(500万円 - 360万円) \times 1/2 = 70万円$
    \item Y社の退職所得は、X社とY社の退職手当等を合計した額からY社の退職所得控除額を控除した額に1/2を乗じて求める。$(500万円+200万円-400万円)\times 1/2 = 150万円$。
  \end{enumerate}
\end{answerbox}

\begin{explanationbox}
  ＜参考＞

タックスアンサー No.2735

\url{https://www.nta.go.jp/taxes/shiraberu/taxanswer/gensen/2735.htm}

\end{explanationbox}


\newpage

\subsection{前年以前に退職手当を受け取っている場合}

\begin{fillblank}[年金1 2025 問題1(6)ウ]
  Ａさんは2025 年 6 月に定年退職し、退職手当として1,000 万円を受け取った。また、2023 年
には60歳到達により個人型確定拠出年金から老齢給付金（一時金）として 500 万円を受け取
っている。

以下の条件において、Ａさんの退職手当に係る退職所得額を計算すると、 （　　A　　）万円で
ある。

＜条件＞
\begin{itemize}
  \item 定年退職した会社での勤続年数は18年2ヵ月
\item 個人型年金加入者期間は10年3ヵ月（全期間で掛金拠出済み）
\item 会社の入社前に個人型年金のみに加入していた期間はなし
\item 勤続年数のうち役員等勤続年数はなし
\item 障害者になったことが退職の直接の原因でない
\end{itemize}
\end{fillblank}

\begin{answerbox}
  10年ルールに抵触しているので、重複期間を勤続年数から控除する。

  控除する勤続年数は月単位を切り捨てるので10年。
  \[
18年2ヶ月 - 10年 = 8年2ヶ月
\]
$\rightarrow$月単位を切り上げて、9年が退職手当にかかる退職所得控除額計算用の勤続年数。
  \begin{flalign*}
退職所得控除額 &= 9年 \times 40万円  \\
          &= 360万円 \\
退職所得 &= (1000万円 - 360万円) \times 1/2 \\
&= 320万円
  \end{flalign*}
\end{answerbox}

\begin{explanationbox}
  10年ルールについては、例えば以下のサイトを参照。

  \url{https://www.freee.co.jp/kb/kb-payroll/retirement-income/?utm_source=google&utm_medium=organic}
\end{explanationbox}

\begin{fillblank}[数理人試験 2024 問題6 設問5]
  勤続25年のＡさんは65歳で定年退職となった。定年退職と同時に確定給付企業年金の加入者
の資格を喪失し、老齢給付金（一時金）を1,500万円受け取った。当該給付の受け取りに際
し、適用される退職所得控除の金額を算定せよ。

＜Ａさんに係る前提＞
\begin{itemize}
  \item 確定給付企業年金および（企業型）確定拠出年金が適用されている。
\item Ａさんが勤務する企業の定年年齢は60歳であったが、Ａさんが60歳を迎える前に定年年齢を65歳
に変更した。また、その際に、確定給付企業年金の資格喪失年齢を65歳に変更した。
\item 40歳から65歳まで継続して勤務しており、定年退職時の確定給付企業年金の加入者期間は25年で
ある。
\item （企業型）確定拠出年金の資格喪失年齢は60歳であり、資格喪失時の加入者期間は20年である。
\item 60歳で（企業型）確定拠出年金の老齢給付金（一時金）を500万円受け取った。
\end{itemize}

\end{fillblank}

\begin{answerbox}
  10年ルールに抵触しているので、重複期間を勤続年数から控除する。

  重複期間を控除すると、勤続期間は$25年-20年=5年$。

  よって退職所得控除額は、
    \begin{flalign*}
退職所得控除額 &= 5年 \times 40万円  \\
          &= 200万円 \\
  \end{flalign*}
  
\end{answerbox}

\begin{explanationbox}
  模範解答は10年ルールの改正前
\end{explanationbox}

\newpage

\subsection{定年延長した場合}

\begin{shortanswer}[数理人試験 2020 問題7 設問1]
  退職一時金にかかる所得税の課税方式・計算方法について簡記せよ。（税率の数値については
言及しなくてよい。）
\end{shortanswer}


\begin{answerbox}
  退職一時金は「退職所得」として扱われ、所得税の計算はいわゆる分離課税の方式が採用されており、他の所得と課税所得を合算せずに、退職所得にかかる所得金額のみで所得税率を決定し、所得税額を計算する。

課税所得は、
\[
(退職一時金の額 - 退職所得控除額) × 1/2
\] で計算され、
退職所得控除額は、20年以下の勤続年数の場合、
\[
40万円 × 勤続年数（ただし80万円未満の場合は、80万円）
\]

20年以上の勤続年数の場合は、
\[
800万円 + 70万円 × (勤続年数 - 20年)
\]

によって計算される。

なお、障害者となったことが原因で退職した場合は、上記の退職所得控除額に100万円加算した額となる。また、前年以前に退職金を受け取ったことがある場合などは、計算が異なることがある。
\end{answerbox}

\begin{shortanswer}[数理人試験 2020 問題7 設問2]
年を延長した場合に旧定年において支払われる退職一時金（確定給付企業年金から支払われ
る脱退一時金を含む）の税務上の所得区分について簡記せよ。なお、旧定年より後でさらに退
職一時金が支給されることはないものとする。

※最終的には税務署の判断に委ねられる事項であるが、例えば定年延長を検討している顧客へ
の一般的な情報提供として、どのようなことを伝えるのが適切かという観点で記述せよ。
\end{shortanswer}

\begin{answerbox}
  定年延長前に入社した従業員については、その支払いをすることについて、相応の理由があると認められる場合は退職所得となる。

定年延長後に入社した従業員については、上記が適用されない可能性が高く、その場合、旧定年で支払われる退職一時金については、会社から支払われる額については給与所得、DBから支払われる額については一時所得となる。
\end{answerbox}

\begin{explanationbox}
  参考：所得税基本通達30-2(5)

  \url{https://www.nta.go.jp/law/tsutatsu/kihon/shotoku/04/04.htm}
\end{explanationbox}

\newpage

\subsection{その他}

\begin{truefalse}[年金1 2018 問題2(3)\textcircled{1}]
  \textbf{以下の文章の正誤を判定し、誤っている場合は正しい内容に改めなさい。}
\\  

  企業型年金の老齢給付金を一時金として受け取るときは退職所得の扱いとなり、退職所得
控除額の計算において勤続年数とみなされるのは、通算加入者等期間である。
\end{truefalse}

\begin{answerbox}
  \textbf{誤}

  勤続年数とみなされるのは、事業主掛金の拠出を伴う企業型年金加入者期間である。
\end{answerbox}

\begin{fillblank}[数理人試験 2025 問題6 設問4]
  次は、所得税法基本通達31-1の記述である。A～Bの空欄に入る語句をそれぞれ記載せよ。
\\

31－1 法第31条第3号に規定する「加入者の退職により支払われるものその他これに類する一時金と
して政令で定めるもの」又は令第72条第2項に規定する「加入員の退職に基因して支払われるもの」
には、確定給付企業年金法の規定に基づいて支払われる退職一時金、公的年金制度の健全性及び信頼性
の確保のための厚生年金保険法等の一部を改正する法律（平成25年法律第63号。以下「平成25年厚
生年金等改正法」という。）第1条((厚生年金保険法の一部改正))の規定による改正前の厚生年金保険
法第9章((厚生年金基金及び企業年金連合会))の規定に基づいて支払われる退職一時金、法人税法附則
第20条第3項((退職年金等積立金に対する法人税の特例))に規定する適格退職年金契約に基づいて支
払われる退職一時金、平成25年厚生年金等改正法附則の規定に基づいて支払われる退職一時金、平成
25 年厚生年金等改正法第 2 条((確定給付企業年金法の一部改正))の規定による改正前の確定給付企業
年金法の規定に基づいて支払われる退職一時金又は確定拠出年金法の規定に基づいて老齢給付金とし
て支払われる一時金のうち、次に掲げる一時金がそれぞれ含まれるものとする。

（1） 確定給付企業年金規約、厚生年金基金規約又は適格退職年金契約に基づいて支給される年金の受
給資格者に対し当該年金に代えて支払われる一時金のうち、退職の日以後当該（　　A　　）までの間
に支払われるもの（（　　A　　）後に支払われる一時金のうち、将来の年金給付の（　　B　　）に代
えて支払われるものを含む。）
（以下、略）
\end{fillblank}

\begin{answerbox}
  \begin{enumerate}[A:]
    \item 年金の受給開始日
    \item 総額
  \end{enumerate}
\end{answerbox}

\newpage